\subsection{Педагогічні умови та стратегія впровадження}\label{subsec:3.3}

Реалізація структурно-функціональної моделі (п.~\ref{subsec:3.1}) та рекомендацій щодо трансформації дисциплін (п.~\ref{subsec:3.2}) потребує визначення педагогічних умов~--- сукупності обставин, що забезпечують ефективність формування ІКК. Педагогічні умови є зовнішніми по відношенню до суб'єкта навчання факторами, які цілеспрямовано створюються в освітньому процесі для забезпечення досягнення визначених результатів.

У цьому підрозділі обґрунтовано 4 педагогічні умови формування ІКК майбутніх дизайнерів, розроблено стратегію поетапного впровадження та запропоновано дизайн експертної оцінки для валідації моделі.

\subsubsection{Обґрунтування педагогічних умов}\label{sssec:3.3.1}

Визначення педагогічних умов ґрунтується на тріангуляції трьох джерел: (а)~емпіричних даних аналізу (розділ~\ref{sec:2}); (б)~теоретичної рамки (п.~\ref{subsec:1.4}); (в)~міжнародного досвіду (п.~\ref{subsec:2.3}). Кожна умова є відповіддю на конкретну проблему, виявлену у ході дослідження.

\textbf{Перша педагогічна умова: системна інтеграція ШІ-інструментів наскрізно через увесь навчальний план.}

\textit{Проблема.} Загальнонаціональний аналіз 122 програм (п.~\ref{subsec:2.1}) виявив 0\% покриття компонентів <<Основи ШІ>>, <<Генеративний ШІ>> та <<Комп'ютерний зір>>. Це не ізольована прогалина окремої дисципліни~--- це \textit{системна відсутність} ШІ-тематики на рівні всієї спеціальності 022~Дизайн в Україні. Кейс-стаді КДПУ (п.~\ref{subsec:2.2}) підтвердив: жоден з 132 проаналізованих силабусів не містить термінів <<штучний інтелект>>, <<генеративний>>, <<нейронна мережа>> чи <<промпт>>. Водночас аналіз виявив нисхідний тренд ІКК-змісту (з 122 до 75--82 ІКК-ключових слів у ОПП між 2018 та 2021~рр.), що свідчить про розмивання, а не посилення ІКК-спрямованого змісту при реформуванні програм.

\textit{Теоретичне обґрунтування.} Модель TPACK~\cite{Mishra2006} наголошує на необхідності інтеграції технологічних знань не як ізольованого компоненту, а як елементу, що пронизує всі аспекти педагогічного процесу. Центральний конструкт TPACK~--- перетин технологічних, педагогічних та змістових знань~--- може бути сформований лише за умови наскрізної інтеграції, коли ШІ-інструменти використовуються у контексті конкретних дизайнерських дисциплін, а не у відриві від них.

\textit{Міжнародний досвід.} Бенчмаркінг 7 провідних програм (п.~\ref{subsec:2.3}) засвідчив: найуспішніші програми (Royal College of Art, Aalto University) інтегрують ШІ-тематику наскрізно через усі дисципліни, а не лише через окремі курси. Цей підхід є принциповим: навіть програми з меншою кількістю спеціалізованих ШІ-дисциплін (Politecnico di Milano~--- 1 ШІ-курс) забезпечують ефективне формування ІКК через наскрізну інтеграцію ШІ-елементів у існуючі курси.

\textit{Зміст умови.} Системна інтеграція передбачає:
\begin{enumerate}
    \item визначення рівня ШІ-інтеграції (базовий, інструментальний або трансформаційний) для кожної дисципліни навчального плану (таблиці~\ref{tab:year1}--\ref{tab:year4});
    \item включення ШІ-компоненту у робочі програми та силабуси відповідних дисциплін з конкретизацією тем, годин та очікуваних результатів;
    \item створення <<ШІ-спіни>> (AI spine)~--- послідовності з 4 спеціалізованих ШІ-дисциплін (таблиця~\ref{tab:new-ai-courses}), що утворюють каркас ШІ-підготовки;
    \item забезпечення \textit{прогресивного зростання} складності ШІ-інтеграції від 1-го до 4-го курсу;
    \item координацію ШІ-компонентів між дисциплінами для уникнення дублювання та забезпечення послідовності.
\end{enumerate}

\textit{Індикатори реалізації:}
\begin{itemize}
    \item частка дисциплін з визначеним рівнем ШІ-інтеграції $\geq$ 70\% (у запропонованому плані~--- 76,7\%);
    \item наявність ШІ-компоненту у силабусах відповідних дисциплін;
    \item наявність <<ШІ-спіни>> з $\geq$ 3 спеціалізованих ШІ-дисциплін;
    \item збільшення кількості ІКК-ключових слів у ОПП порівняно з поточним станом.
\end{itemize}

\textbf{Друга педагогічна умова: трансформація студійної педагогіки від моделі <<дизайнер-виробник>> до моделі <<дизайнер-виробник + дизайнер-куратор>>.}

\textit{Проблема.} Оглядове дослідження 156 наукових праць~\cite{Musiienko2026ScopingReview} зафіксувало парадигмальний зсув у дизайн-освіті: генеративний ШІ трансформує роль дизайнера від виключно виробника (maker), що створює артефакти своїми руками, до куратора (curator), що оцінює, обирає та редагує згенеровані ШІ варіанти. Цей зсув потребує відповідної трансформації студійної педагогіки~--- <<підписової педагогіки>> дизайнерських професій~\cite{Shulman2005}.

Поточний стан студійної педагогіки в КДПУ (та в Україні загалом, за даними п.~\ref{subsec:2.1}) орієнтований виключно на модель виробника: студент створює артефакти вручну або за допомогою традиційних цифрових інструментів, а оцінюється якість виконання. Модель куратора~--- здатність критично оцінювати, відбирати та редагувати ШІ-генеровані варіанти~--- не формується, оскільки ШІ-інструменти відсутні у навчальному процесі.

\textit{Теоретичне обґрунтування.} Теорія рефлексивної практики~\cite{Schon1983} обґрунтовує подвійну модель <<виробник + куратор>>: рефлексія-в-дії (reflection-in-action) активізується як при ручному створенні, так і при курації ШІ-результатів; рефлексія-про-дію (reflection-on-action) збагачується новим виміром~--- порівнянням стратегій ручного та ШІ-базованого створення. Студійна педагогіка, адаптована для подвійної моделі, формує більш розвинений рефлексивний компонент ІКК.

Важливо підкреслити: модель <<виробник + куратор>> не означає відмови від ручних навичок. Ручне створення (рисунок, живопис, макетування) залишається фундаментом дизайнерської підготовки (принцип доповнення, а не заміщення, п.~\ref{subsec:3.1}). Курація ШІ-результатів є \textit{додатковою} компетенцією, що доповнює, а не замінює виробничу.

\textit{Зміст умови.} Трансформація студійної педагогіки передбачає:
\begin{enumerate}
    \item модифікацію студійного циклу: класичний цикл <<ідея $\to$ прототип $\to$ критика $\to$ вдосконалення>> доповнюється етапами ШІ-генерації та порівняльного аналізу (п.~\ref{subsec:3.1});
    \item оновлення формату критики (crit): публічне обговорення включає аналіз стратегії використання ШІ, обґрунтування інструментального вибору та етичну рефлексію;
    \item введення завдань <<порівняння ручне vs. ШІ>> у студійних курсах (п.~\ref{subsec:3.2});
    \item формування компетенцій куратора: відбір серед ШІ-варіантів, оцінка якості генерації, редагування та доопрацювання ШІ-результатів;
    \item збереження обсягу ручної практики~--- завдання з ШІ додаються, а не замінюють ручні завдання.
\end{enumerate}

\textit{Індикатори реалізації:}
\begin{itemize}
    \item наявність завдань <<ручне vs. ШІ>> у $\geq$ 50\% студійних курсів;
    \item включення ШІ-виміру у формат критики;
    \item фіксація компетенцій куратора у результатах навчання відповідних дисциплін;
    \item збереження обсягу годин на ручну практику на рівні $\geq$ 90\% від поточного.
\end{itemize}

\textbf{Третя педагогічна умова: узгодження розвитку компетенцій з вимогами роботодавців та ринку праці.}

\textit{Проблема.} Опитування роботодавців (п.~\ref{subsec:2.2}) виявило чіткий сигнал: роботодавці програми <<Дизайн одягу>> прямо рекомендують <<використання штучного інтелекту у сфері дизайну>>. Оцінка ІКТ-навичок випускників~--- 3,5--4,0 із 5~--- нижче середнього. Серед пріоритетів вдосконалення роботодавці виділяють зміни у змісті навчальних програм, збільшення обсягу практичної підготовки та оновлення матеріально-технічної бази.

Водночас опитування випускників (п.~\ref{subsec:2.2}) засвідчило суттєву незадоволеність: 33,3\% випускників ГД та 15,4\% випускників ДО не задоволені рівнем ІКТ-підготовки. Випускники називають конкретне програмне забезпечення, необхідне на ринку праці (Adobe After Effects, Figma, Blender, Cinema~4D тощо), проте жоден не згадав ШІ-інструменти~--- ймовірно, через їх відсутність у навчальному процесі.

Аналіз висновків експертних груп НАЗЯВО (п.~\ref{subsec:2.2}) додатково підтвердив потребу в оновленні: технологічний компонент (70--80 згадувань) представлений переважно через згадки традиційного обладнання та ПЗ, без ШІ-тематики.

\textit{Теоретичне обґрунтування.} Компетентнісний підхід~\cite{Musiienko2025DesignEducation} передбачає, що результати навчання мають відповідати реальним потребам ринку праці та суспільства. Стандарт В2 Дизайн визначає зв'язок з роботодавцями як один з механізмів забезпечення якості освіти. Аналіз прогалин стандарту В2 (п.~\ref{subsec:1.4}) виявив відсутність ШІ-компетентностей, що підтверджує необхідність випереджального оновлення програм відносно стандарту.

\textit{Зміст умови.} Узгодження з вимогами роботодавців передбачає:
\begin{enumerate}
    \item регулярне опитування роботодавців та випускників (щонайменше раз на 2 роки) щодо актуальних технологій та навичок, з включенням блоку питань про ШІ-інструменти;
    \item залучення практикуючих дизайнерів до навчального процесу як запрошених лекторів (guest lecturers) та менторів проєктів, зокрема дизайнерів, що активно використовують ШІ;
    \item формування баз практики з урахуванням ШІ-використання на підприємствах (п.~\ref{subsec:3.2});
    \item оновлення переліку ШІ-інструментів у навчальному плані щонайменше щорічно відповідно до актуальних технологічних тенденцій;
    \item проведення спільних проєктів <<університет~--- підприємство>> з використанням ШІ-інструментів;
    \item включення до вимог дипломного проєкту обов'язкового ШІ-компоненту (документація використання ШІ, рефлексія).
\end{enumerate}

\textit{Індикатори реалізації:}
\begin{itemize}
    \item проведення опитування роботодавців з ШІ-блоком $\geq$ 1 раз на 2 роки;
    \item $\geq$ 2 запрошених лекторів-практиків з ШІ-досвідом на рік;
    \item $\geq$ 30\% баз практики з підтвердженим ШІ-використанням;
    \item щорічне оновлення переліку ШІ-інструментів у силабусах;
    \item зростання задоволеності випускників ІКТ-підготовкою (цільовий показник: $<$ 15\% незадоволених).
\end{itemize}

\textbf{Четверта педагогічна умова: безперервний моніторинг формування ІКК через портфоліо-орієнтоване оцінювання та рефлексивну практику.}

\textit{Проблема.} Аналіз висновків НАЗЯВО (п.~\ref{subsec:2.2}) виявив, що рефлексивний компонент має найменше представництво (11--12 згадувань порівняно з 70--114 для інших компонентів). Це свідчить про відсутність системи цілеспрямованого формування та моніторингу рефлексивного компоненту ІКК. Водночас рефлексивний компонент є критично важливим в ШІ-еру: здатність критично оцінювати результати ШІ-генерації, обґрунтовувати інструментальний вибір та рефлексувати щодо етичних аспектів~--- це те, що відрізняє компетентного дизайнера від некритичного користувача ШІ.

\textit{Теоретичне обґрунтування.} Рефлексивна практика за Шоном~\cite{Schon1983} є невід'ємною складовою студійної педагогіки: рефлексія-в-дії та рефлексія-про-дію забезпечують усвідомлене опанування професійних навичок. Компетентнісний підхід передбачає безперервне оцінювання (continuous assessment) як альтернативу одноразовим іспитам. Додаткового значення рефлексивний компонент набуває у контексті Закону України <<Про штучний інтелект>>~\cite{ZakonAI2025}, стаття~8(4) якого вимагає розкриття ШІ-генерованого контенту в академічних роботах із зазначенням методології, а стаття~29 кваліфікує подання ШІ-результатів без розкриття як порушення академічної доброчесності. Це створює законодавчу підставу для систематичного формування рефлексивної практики ШІ-використання. Портфоліо як інструмент оцінювання визначено найкращою практикою у міжнародному бенчмаркінгу (п.~\ref{subsec:2.3}): RMIT University використовує цифрове портфоліо, що демонструє як ШІ-проєкти, так і традиційні роботи.

\textit{Зміст умови.} Безперервний моніторинг передбачає:
\begin{enumerate}
    \item впровадження \textit{цифрового портфоліо} як наскрізного інструменту оцінювання ІКК (від 1-го до 4-го курсу). Портфоліо оцінюється за критеріями, визначеними у діагностично-результативному блоці моделі (таблиця~\ref{tab:ikk-criteria});
    \item систематичне ведення \textit{рефлексивного щоденника} у студійних курсах з фіксацією досвіду взаємодії з ШІ-інструментами;
    \item періодичне \textit{самооцінювання} за 4 компонентами ІКК (початок та кінець кожного семестру) з використанням стандартизованого опитувальника;
    \item портфоліо-перегляд (portfolio review) наприкінці кожного навчального року з оцінкою динаміки формування ІКК;
    \item включення ШІ-виміру у публічні критики (crit sessions): обов'язкове обговорення стратегії використання ШІ;
    \item фінальний захист портфоліо як компонент державної атестації: портфоліо з ШІ-проєктами та рефлексивним аналізом ролі ШІ у професійному розвитку.
\end{enumerate}

\textit{Індикатори реалізації:}
\begin{itemize}
    \item 100\% здобувачів ведуть цифрове портфоліо;
    \item проведення портфоліо-перегляду $\geq$ 1 раз на рік;
    \item наявність рефлексивного щоденника у $\geq$ 70\% студійних курсів;
    \item проведення самооцінювання ІКК $\geq$ 2 рази на рік;
    \item включення портфоліо-захисту у процедуру державної атестації.
\end{itemize}

\subsubsection{Взаємозв'язок педагогічних умов}\label{sssec:3.3.2}

Чотири педагогічні умови утворюють цілісну систему, де кожна умова підтримує та посилює інші:

\begin{itemize}
    \item \textbf{Умова~1} (системна ШІ-інтеграція) створює \textit{змістовий каркас}, у межах якого реалізуються всі інші умови. Без наскрізної інтеграції ШІ в навчальний план неможливо трансформувати студійну педагогіку (умова~2), узгодити підготовку з вимогами ринку (умова~3) чи здійснювати моніторинг ІКК (умова~4).

    \item \textbf{Умова~2} (трансформація студійної педагогіки) забезпечує \textit{методичну основу} для ефективного використання ШІ-інструментів (умова~1) у контексті студійного навчання. Модифікована студійна педагогіка також генерує дані (рефлексивні записи, crit-протоколи) для моніторингу ІКК (умова~4).

    \item \textbf{Умова~3} (узгодження з ринком) забезпечує \textit{актуальність} змісту ШІ-інтеграції (умова~1), визначаючи, які саме інструменти та компетенції є затребуваними. Зворотний зв'язок від роботодавців інформує оновлення навчального плану та критеріїв оцінювання (умова~4).

    \item \textbf{Умова~4} (моніторинг через портфоліо) забезпечує \textit{зворотний зв'язок} для всіх інших умов: результати моніторингу виявляють прогалини у ШІ-інтеграції (умова~1), необхідність коригування студійних методів (умова~2) та відповідність очікуванням ринку (умова~3).
\end{itemize}

Цей системний взаємозв'язок відображає структуру моделі: діагностично-результативний блок (умова~4) забезпечує зворотний зв'язок для мотиваційно-цільового (умова~3), змістового (умова~1) та організаційно-технологічного (умова~2) блоків.

\subsubsection{Стратегія поетапного впровадження}\label{sssec:3.3.3}

Впровадження моделі трансформації потребує поетапного підходу, що враховує реальні обмеження: інституційну інерцію, потребу у підвищенні кваліфікації викладачів, обмеженість ресурсів (зокрема, зауважену студентами ГД проблему застарілого обладнання та ПЗ, п.~\ref{subsec:2.2}). Стратегія передбачає 4 етапи, що відповідають рокам навчання, але реалізуються послідовно на інституційному рівні.

\textbf{Етап 1. Усвідомлення та підготовка (Рік впровадження 1).}

\textit{Мета}: створення передумов для системної трансформації.

\textit{Заходи:}
\begin{enumerate}
    \item Проведення серії семінарів для викладачів з ШІ-інструментами для дизайну (Midjourney, Adobe Firefly, Figma AI, ChatGPT). Цільова аудиторія: усі викладачі кафедри дизайну (10--15 осіб). Формат: 4 семінари по 4 години = 16 годин.
    \item Введення дисципліни <<Цифрова грамотність та основи AI>> (4 кредити, семестр~1) для студентів 1-го курсу.
    \item Пілотна інтеграція ШІ-елементів у 3--4 існуючі дисципліни (Композиція, Комп'ютерна графіка, Графічний дизайн, Дизайн-мислення) на базовому рівні.
    \item Розробка шаблону рефлексивного щоденника та запуск пілотного портфоліо.
    \item Проведення опитування роботодавців з блоком питань про ШІ.
\end{enumerate}

\textit{Очікувані результати}: 100\% викладачів ознайомлені з базовими ШІ-інструментами; 1 нова дисципліна введена; 3--4 існуючі дисципліни оновлені на базовому рівні; базові дані для моніторингу зібрано.

\textbf{Етап 2. Інструментальна інтеграція (Рік впровадження 2).}

\textit{Мета}: системне впровадження ШІ-інструментів у фахові дисципліни.

\textit{Заходи:}
\begin{enumerate}
    \item Оновлення силабусів 8--10 фахових дисциплін з включенням ШІ-модулів інструментального рівня (таблиці~\ref{tab:digital-transform}--\ref{tab:professional-transform}).
    \item Модифікація формату критики (crit) з включенням ШІ-виміру.
    \item Введення завдань <<ручне vs. ШІ>> у студійні курси.
    \item Формування бази практики з ШІ-використанням.
    \item Перший портфоліо-перегляд наприкінці року.
    \item Поглиблене підвищення кваліфікації викладачів: промпт-інженерія, інтеграція ШІ у конкретні дисципліни. Формат: індивідуальні консультації + 2 майстер-класи.
\end{enumerate}

\textit{Очікувані результати}: 8--10 дисциплін мають визначений рівень ШІ-інтеграції; формат критики модифіковано; база практики з ШІ $\geq$ 20\%; перший цикл моніторингу завершено.

\textbf{Етап 3. Поглиблена інтеграція (Рік впровадження 3).}

\textit{Мета}: введення спеціалізованих ШІ-дисциплін та досягнення трансформаційного рівня.

\textit{Заходи:}
\begin{enumerate}
    \item Введення дисциплін <<Генеративний дизайн та AI>> (6 кредитів) та <<Взаємодія людина-AI в дизайні>> (4 кредити).
    \item Трансформаційний рівень ШІ-інтеграції у виробничій практиці.
    \item Повномасштабне впровадження портфоліо-оцінювання.
    \item Залучення запрошених лекторів-практиків з ШІ-досвідом ($\geq$ 2 на рік).
    \item Створення ШІ-лабораторії або виділеного простору з обладнанням для роботи з ШІ (GPU, ліцензії).
    \item Другий цикл опитування роботодавців; коригування програми за результатами.
\end{enumerate}

\textit{Очікувані результати}: 2 нові ШІ-дисципліни введені; трансформаційний рівень досягнуто у 4+ дисциплінах; база практики з ШІ $\geq$ 30\%; ШІ-лабораторія функціонує.

\textbf{Етап 4. Повна реалізація та оцінка (Рік впровадження 4).}

\textit{Мета}: завершення трансформації, перший повний випуск за оновленою програмою, оцінка результатів.

\textit{Заходи:}
\begin{enumerate}
    \item Введення дисципліни <<Критичний дизайн та етика AI>> (4 кредити).
    \item Включення ШІ-компоненту у вимоги до дипломного проєкту.
    \item Включення портфоліо-захисту у процедуру державної атестації.
    \item Проведення комплексної оцінки: опитування випускників, роботодавців, порівняння з базовими показниками (п.~\ref{subsec:2.2}).
    \item Публікація результатів трансформації та рекомендацій.
    \item Коригування моделі за результатами першого повного циклу.
\end{enumerate}

\textit{Очікувані результати}: повний 4-річний цикл реалізовано; 4 нові ШІ-дисципліни функціонують; $>$ 75\% дисциплін мають ШІ-інтеграцію; перший випуск за оновленою програмою; дані для порівняння <<до>> та <<після>> трансформації зібрано.

\subsubsection{Хронологічна матриця впровадження}\label{sssec:3.3.4}

\begin{longtable}{@{}p{0.32\textwidth}cccc@{}}
\caption{Матриця поетапного впровадження}\label{tab:implementation-matrix} \\
\hline
\textbf{Захід} & \textbf{Етап 1} & \textbf{Етап 2} & \textbf{Етап 3} & \textbf{Етап 4} \\
\hline
\endfirsthead
\hline
\textbf{Захід} & \textbf{Етап 1} & \textbf{Етап 2} & \textbf{Етап 3} & \textbf{Етап 4} \\
\hline
\endhead
Семінари для викладачів & $\bullet$ & $\bullet$ & & \\
\addlinespace
Нова дисципліна: Цифрова грамотність та основи AI & $\bullet$ & $\bullet$ & $\bullet$ & $\bullet$ \\
\addlinespace
Пілотна ШІ-інтеграція (3--4 дисципліни) & $\bullet$ & & & \\
\addlinespace
Оновлення силабусів (8--10 дисциплін) & & $\bullet$ & $\bullet$ & $\bullet$ \\
\addlinespace
Модифікація формату критики & & $\bullet$ & $\bullet$ & $\bullet$ \\
\addlinespace
Нові дисципліни: Генеративний дизайн та AI; Взаємодія людина-AI & & & $\bullet$ & $\bullet$ \\
\addlinespace
Нова дисципліна: Критичний дизайн та етика AI & & & & $\bullet$ \\
\addlinespace
Портфоліо-оцінювання & пілот & часткове & повне & повне \\
\addlinespace
Рефлексивний щоденник & пілот & 50\% курсів & 70\% курсів & 70\%+ курсів \\
\addlinespace
Опитування роботодавців (з ШІ-блоком) & $\bullet$ & & $\bullet$ & \\
\addlinespace
Запрошені лектори-практики & & $\geq$1 & $\geq$2 & $\geq$2 \\
\addlinespace
ШІ-лабораторія & планування & закупівля & запуск & функціонування \\
\addlinespace
ШІ у дипломному проєкті & & & рекомендовано & обов'язково \\
\addlinespace
Комплексна оцінка результатів & базові дані & проміжна & проміжна & фінальна \\
\hline
\end{longtable}

\subsubsection{Ресурсне забезпечення впровадження}\label{sssec:3.3.5}

Реалізація моделі потребує ресурсного забезпечення за кількома напрямками:

\textbf{1. Матеріально-технічне забезпечення.} Опитування студентів (п.~\ref{subsec:2.2}) виявило статистично значущу незадоволеність матеріально-технічним забезпеченням ($H = 8{,}708$, $p = 0{,}013$), особливо у програмі ГД. Студенти згадують <<застаріле програмне забезпечення (Photoshop 2003)>> та <<старі комп'ютери>>. Мінімальні вимоги для ШІ-інтеграції:
\begin{itemize}
    \item оновлення комп'ютерного парку: $\geq$ 15 робочих станцій з сучасними GPU (NVIDIA RTX 3060+) для запуску локальних ШІ-моделей (Stable Diffusion);
    \item придбання ліцензій на хмарні ШІ-сервіси: Midjourney, Adobe Creative Cloud (з Firefly), Figma Professional;
    \item стабільне Інтернет-з'єднання з достатньою пропускною здатністю для хмарних ШІ-сервісів;
    \item виділений простір (ШІ-лабораторія) з обладнаними робочими місцями.
\end{itemize}

\textbf{2. Кадрове забезпечення.} Підвищення кваліфікації викладачів є критичною передумовою впровадження. Програма підвищення кваліфікації передбачає:
\begin{itemize}
    \item базовий модуль (16 годин): ознайомлення з ШІ-інструментами для дизайну~--- для всіх викладачів кафедри;
    \item поглиблений модуль (32 години): промпт-інженерія, інтеграція ШІ у конкретну дисципліну, портфоліо-оцінювання~--- для викладачів фахових дисциплін;
    \item стажування на підприємствах, що використовують ШІ~--- для 2--3 викладачів на рік;
    \item можливе залучення зовнішніх фахівців (дизайнерів-практиків з ШІ-досвідом) для викладання нових дисциплін на початковому етапі.
\end{itemize}

\textbf{3. Навчально-методичне забезпечення.} Розробка нових навчально-методичних матеріалів:
\begin{itemize}
    \item робочі програми та силабуси 4 нових ШІ-дисциплін;
    \item оновлені силабуси існуючих дисциплін з ШІ-модулями;
    \item методичні рекомендації для викладачів щодо інтеграції ШІ у конкретні типи дисциплін;
    \item шаблони рефлексивного щоденника та портфоліо;
    \item опитувальник для самооцінки ІКК;
    \item банк навчальних завдань з ШІ-компонентом.
\end{itemize}

\subsubsection{Експертна оцінка моделі}\label{sssec:3.3.6}

Для валідації запропонованої моделі розроблено дизайн експертної оцінки, що передбачає залучення 5--7 експертів~--- фахівців у галузі дизайн-освіти, ІКТ в освіті та/або практикуючих дизайнерів з досвідом використання ШІ.

\textbf{Критерії відбору експертів:}
\begin{itemize}
    \item науковий ступінь (PhD або Dr. habil.) та/або $\geq$ 10 років досвіду у дизайн-освіті;
    \item досвід розробки або рецензування освітніх програм спеціальності 022~Дизайн;
    \item обізнаність з ШІ-інструментами для дизайну (бажано);
    \item досвід експертизи НАЗЯВО (бажано).
\end{itemize}

\textbf{Опитувальник експертної оцінки.} Опитувальник складається з 5 блоків, що відповідають блокам структурно-функціональної моделі, та загального блоку. Кожне твердження оцінюється за 5-бальною шкалою (1~--- повністю не погоджуюсь, 5~--- повністю погоджуюсь).

\textit{Блок~1. Мотиваційно-цільовий блок} (5 тверджень):
\begin{enumerate}
    \item Мета моделі (формування ІКК через системну трансформацію навчальних програм з інтеграцією ШІ) є актуальною та обґрунтованою.
    \item Завдання моделі за 4 компонентами ІКК є конкретними та досяжними.
    \item Соціальне замовлення (дані опитувань роботодавців, випускників, аналіз програм) достатньо обґрунтовує необхідність трансформації.
    \item Формулювання компонентів ІКК (інформаційно-аналітичний, комунікативний, технологічний, рефлексивний) є повним та адекватним для дизайнерської професії.
    \item Модель враховує сучасні вимоги ринку праці до дизайнерів.
\end{enumerate}

\textit{Блок~2. Теоретико-методологічний блок} (4 твердження):
\begin{enumerate}
    \setcounter{enumi}{5}
    \item Теоретичні підходи (TPACK, студійна педагогіка, компетентнісний підхід) є обґрунтованим підґрунтям моделі.
    \item Дидактичні принципи (системність, доповнення, прогресивність, практикоорієнтованість, рефлексивність, відповідність ринку) є повними та послідовними.
    \item Принцип <<доповнення, а не заміщення>> (збереження ручних навичок) є необхідним для дизайн-освіти.
    \item Трансформація студійної педагогіки від <<виробника>> до <<виробника + куратора>> є педагогічно обґрунтованою.
\end{enumerate}

\textit{Блок~3. Змістовий блок} (4 твердження):
\begin{enumerate}
    \setcounter{enumi}{9}
    \item Відображення 4 компонентів ІКК на навчальні дисципліни є коректним та повним.
    \item Рекомендації щодо трансформації студійних курсів (Живопис, Композиція) є реалістичними та ефективними.
    \item Рекомендації щодо трансформації цифрових та професійних курсів є реалістичними та ефективними.
    \item 240-кредитний навчальний план з <<ШІ-спіною>> є збалансованим та реалізовуваним.
\end{enumerate}

\textit{Блок~4. Організаційно-технологічний блок} (4 твердження):
\begin{enumerate}
    \setcounter{enumi}{13}
    \item Набір ШІ-інструментів (Midjourney, DALL-E, Stable Diffusion, Adobe Firefly, Figma AI, RunwayML) є актуальним та достатнім.
    \item Методи навчання (модифікована студія, проєктне навчання з ШІ, критика з ШІ-виміром) є ефективними для формування ІКК.
    \item Поетапна стратегія впровадження (4 етапи) є реалістичною.
    \item Ресурсні вимоги (обладнання, кваліфікація викладачів, методичне забезпечення) є адекватними.
\end{enumerate}

\textit{Блок~5. Діагностично-результативний блок} (3 твердження):
\begin{enumerate}
    \setcounter{enumi}{17}
    \item Критерії оцінювання ІКК за 4 компонентами є повними та вимірюваними.
    \item Портфоліо-орієнтоване оцінювання є найбільш адекватним інструментом для оцінки ІКК дизайнерів.
    \item 4-рівнева шкала сформованості ІКК (початковий, базовий, достатній, високий) є діагностичною.
\end{enumerate}

\textit{Блок~6. Загальна оцінка} (3 твердження):
\begin{enumerate}
    \setcounter{enumi}{20}
    \item Модель загалом є цілісною та внутрішньо узгодженою.
    \item Модель може бути адаптована для інших програм спеціальності 022~Дизайн в Україні.
    \item Реалізація моделі суттєво підвищить рівень сформованості ІКК випускників.
\end{enumerate}

Кожен блок також містить відкрите питання: <<Ваші зауваження та рекомендації щодо цього блоку>>.

\textbf{Методика обробки результатів.} Для кожного твердження обчислюються: середнє значення ($M$), стандартне відхилення ($SD$), мінімальне та максимальне значення. Тоді обраховується коефіцієнт конкордації Кендалла ($W$) для оцінки узгодженості думок експертів. Прийнятним вважається $W \geq 0{,}50$ при $p < 0{,}05$. Твердження з середнім значенням $M < 3{,}5$ потребують перегляду; з $M \geq 4{,}0$~--- вважаються схваленими; $3{,}5 \leq M < 4{,}0$~--- потребують уточнення.

Відкриті відповіді аналізуються методом тематичного кодування для виявлення системних зауважень та рекомендацій.

\subsubsection{Обмеження та перспективи}\label{sssec:3.3.7}

Запропонована модель має кілька обмежень, що визначають напрямки подальших досліджень:

\begin{enumerate}
    \item \textbf{Відсутність експериментальної верифікації}. Модель розроблена на основі аналізу даних та теоретичного обґрунтування, проте не пройшла експериментальну перевірку (педагогічний експеримент з контрольною та експериментальною групами). Це є обмеженням дослідження та перспективним напрямком подальшої роботи. Водночас підхід, заснований на аналізі даних (data-driven), є обґрунтованою альтернативою в ситуації, коли трансформація ще не впроваджена.

    \item \textbf{Швидка зміна ШІ-інструментів}. Конкретні інструменти (Midjourney, DALL-E, Stable Diffusion тощо) можуть змінитися або зникнути протягом 2--3 років. Модель адресує це обмеження через орієнтацію на \textit{трансферні навички} (промпт-інженерія, критична оцінка, інтеграція ШІ у процес), а не на конкретні продукти. Щорічне оновлення переліку інструментів є частиною третьої педагогічної умови.

    \item \textbf{Контекстна специфічність}. Модель розроблена для програм спеціальності 022~Дизайн, зокрема з урахуванням специфіки КДПУ. Адаптація для інших спеціальностей (архітектура, образотворче мистецтво, кіно) потребує додаткового дослідження. Водночас загальна структура (5 блоків, 4 компоненти ІКК, 4 педагогічні умови) може бути використана як базова модель.

    \item \textbf{Ресурсні обмеження}. Реалізація моделі потребує значних ресурсів (обладнання, ліцензії, підвищення кваліфікації). У контексті українських реалій (обмежене фінансування, наслідки війни) повна реалізація може потребувати більше часу, ніж передбачені 4 роки. Поетапна стратегія впровадження (п.~\ref{subsec:3.3}) частково адресує це обмеження.

    \item \textbf{Етичні питання ШІ в освіті}. Інтеграція ШІ у навчальний процес породжує етичні питання: академічна доброчесність (використання ШІ для виконання завдань), конфіденційність даних (завантаження студентських робіт до хмарних сервісів), нерівність доступу (платні ШІ-сервіси). Четверта педагогічна умова (рефлексивна практика) та дисципліна <<Критичний дизайн та етика AI>> частково адресують ці питання, проте вони потребують постійної уваги.
\end{enumerate}

Попри зазначені обмеження, запропонована модель є першою спробою системної трансформації навчальних програм дизайну в Україні для формування ІКК з урахуванням ШІ-виміру. Вона ґрунтується на масштабній емпіричній базі (122 програми, 132 силабуси, опитування випускників, роботодавців, студентів, висновки НАЗЯВО, міжнародний бенчмаркінг) та інтегративній теоретичній рамці (TPACK, студійна педагогіка, компетентнісний підхід), що забезпечує її наукову обґрунтованість.
